% Part: first-order-logic
% Chapter: completeness
% Section: complete-sets

% Definition of complete consistent sets. Properties of complete sets required
% for completeness proved are in provability.tex in the
% chapter on the proof system used.

\documentclass[../../../include/open-logic-section]{subfiles}

\begin{document}

\iftag{FOL}
      {\olfileid{fol}{com}{ccs}}
      {\olfileid{pl}{com}{ccs}}
      
\olsection{Himpunan \usetoken{P}{sentence} kang Jangkep lan Konsisten}

\begin{defn}[Himpunan jangkep]
\ollabel{def:complete-set} Himpunan~$\Gamma$ saka !!{sentence}s iku
\emph{!!{complete}} yen lan mung yen kanggo sembarang !!{sentence}~$!A$,
$!A \in \Gamma$ utawa $\lnot !A \in \Gamma$.
\end{defn}

\begin{explain}
Himpunan ukara kang !!{complete} ora ninggal pitakon kang durung kajawab.
Kanggo sembarang !!{sentence}~$!A$, $\Gamma$ ``ngandharake'' apa $!A$
bener utawa luput. Wigatine himpunan !!{complete} ora mung tumrap bukti
teorema kelengkapan. Umpamane, teori kang !!{complete} lan bisa
diaksiomatisasi mesthi bisa diputusake.
\end{explain}

\begin{explain}
Himpunan konsisten kang !!{complete} wigati ing bukti kelengkapan amarga
kita bisa njamin yen saben himpunan konsisten saka
!!{sentence}s~$\Gamma$ kamot ing himpunan konsisten kang !!a{complete}
$\Gamma^*$. Himpunan konsisten kang !!a{complete} ngemot, kanggo saben
!!{sentence}~$!A$, $!A$ utawa negasine $\lnot !A$, nanging ora
kalorone. Iki mligine bener kanggo \iftag{FOL}{!!{sentence}s
atomik}{variabel proposisional}, mula saka himpunan
konsisten kang !!a{complete}\iftag{FOL}{ ing basa kang wis digedhekake kanthi trep nganggo
!!{constant}s}{}, kita bisa mbangun
\iftag{FOL}{!!a{structure}}{!!a{valuation}} kang
\iftag{FOL}{interpretasi !!{predicate}s}{nilai bebener kang dipasangake
marang variabel proposisional} ditemtokake miturut
\iftag{FOL}{!!{sentence}s atomik}{!!{propositional variable}s} endi kang
ana ing~$\Gamma^*$. Banjur bisa dibuktekake yen
\iftag{FOL}{!!{structure}}{!!{valuation}} iki ndadekake kabeh
!!{sentence}s ing~$\Gamma^*$ (lan mula uga kabeh kang ana
ing~$\Gamma$) bener. Bukti kasunyatan pungkasan iki mbutuhake yen
$\lnot !A \in \Gamma^*$ yen lan mung yen $!A \notin \Gamma^*$,
$(!A \lor !B) \in \Gamma^*$ yen lan mung yen $!A \in \Gamma^*$ utawa
$!B \in \Gamma^*$, lan sateruse.
\end{explain}

Ing sabanjure, kita bakal kerep nggunakake tanpa nyebutake
sipat refleksivitas, monotonisitas, lan transitivitas saka $\Proves$
(delengen
\iftag{FOL}{%
  \tagrefs{prfSC/{fol:seq:ptn:sec},prfND/{fol:ntd:ptn:sec},prfAX/{fol:axd:ptn:sec},prfTab/{fol:tab:ptn:sec}}}{%
  \tagrefs{prfSC/{pl:seq:ptn:sec},prfND/{pl:ntd:ptn:sec},prfAX/{pl:axd:ptn:sec},prfTab/{pl:tab:ptn:sec}}}).

\begin{prop}
\ollabel{prop:ccs}
Upamana $\Gamma$ iku !!{complete} lan konsisten. Banjur:
\begin{enumerate}
\item \ollabel{prop:ccs-prov-in} Yen $\Gamma \Proves !A$, mula $!A \in
  \Gamma$.

\tagitem{prvAnd}{\ollabel{prop:ccs-and} $!A \land !B \in \Gamma$
  yen lan mung yen $!A \in \Gamma$ lan $!B \in \Gamma$ kalorone.}{}

\tagitem{prvOr}{\ollabel{prop:ccs-or} $!A \lor !B \in \Gamma$ yen lan
  mung yen $!A \in \Gamma$ utawa $!B \in \Gamma$.}{}

\tagitem{prvIf}{\ollabel{prop:ccs-if} $!A \lif !B \in \Gamma$ yen lan
  mung yen $!A \notin \Gamma$ utawa $!B \in \Gamma$.}{}
\end{enumerate}
\end{prop}

\begin{proof}
Kanggo kabeh bagean ing ngisor iki, upamana $\Gamma$ iku !!{complete}
lan konsisten.
\begin{enumerate}
\item Yen $\Gamma \Proves !A$, mula $!A \in \Gamma$.

Upamana $\Gamma \Proves !A$. Kanggo nggayuh cengkah, upamana $!A
\notin \Gamma$. Amarga $\Gamma$ iku !!{complete}, $\lnot !A \in \Gamma$.
Miturut
\iftag{FOL}{%
  \tagrefs{prfSC/{fol:seq:prv:prop:explicit-inc},
    prfND/{fol:ntd:prv:prop:explicit-inc},
    prfAX/{fol:axd:prv:prop:explicit-inc},
    prfTab/{fol:tab:prv:prop:explicit-inc}}}{%
  \tagrefs{prfSC/{pl:seq:prv:prop:explicit-inc},
    prfND/{pl:ntd:prv:prop:explicit-inc},
    prfAX/{pl:axd:prv:prop:explicit-inc},
    prfTab/{pl:tab:prv:prop:explicit-inc}}},
$\Gamma$ ora konsisten. Iki cengkah karo hipotesis yen $\Gamma$
konsisten. Mula, ora mungkin $!A \notin \Gamma$, dadi $!A \in
\Gamma$.

\tagitem{defAnd}{}{%
\iftag{probAnd}{Gladhen.}{%
$!A \land !B \in \Gamma$ yen lan mung yen $!A \in \Gamma$ lan
$!B \in \Gamma$ kalorone:

Kanggo arah maju, upamana $!A \land !B \in \Gamma$. Banjur miturut
\iftag{FOL}{%
  \tagrefs{prfSC/{fol:seq:ppr:prop:provability-land},
    prfND/{fol:ntd:ppr:prop:provability-land},
    prfAX/{fol:axd:ppr:prop:provability-land},
    prfTab/{fol:tab:ppr:prop:provability-land}}}{%
  \tagrefs{prfSC/{pl:seq:ppr:prop:provability-land},
    prfND/{pl:ntd:ppr:prop:provability-land},
    prfAX/{pl:axd:ppr:prop:provability-land},
    prfTab/{pl:tab:ppr:prop:provability-land}}}, butir~(1),
$\Gamma \Proves !A$ lan $\Gamma \Proves !B$. Miturut
\olref{prop:ccs-prov-in}, $!A \in \Gamma$ lan $!B \in \Gamma$,
kaya kang dibutuhake.

Kanggo arah kosok baline, upamana $!A \in \Gamma$ lan $!B \in
\Gamma$. Miturut
\iftag{FOL}{%
  \tagrefs{prfSC/{fol:seq:ppr:prop:provability-land},
    prfND/{fol:ntd:ppr:prop:provability-land},
    prfAX/{fol:axd:ppr:prop:provability-land},
    prfTab/{fol:tab:ppr:prop:provability-land}}}{%
  \tagrefs{prfSC/{pl:seq:ppr:prop:provability-land},
    prfND/{pl:ntd:ppr:prop:provability-land},
    prfAX/{pl:axd:ppr:prop:provability-land},
    prfTab/{pl:tab:ppr:prop:provability-land}}}, butir~(2),
$\Gamma \Proves !A \land !B$. Miturut \olref{prop:ccs-prov-in},
$!A \land !B \in \Gamma$.}}

\tagitem{defOr}{}{%
\iftag{probOr}{Gladhen.}{%
Kaping pisan, kita nuduhake yen $!A \lor !B \in \Gamma$, mula
$!A \in \Gamma$ utawa $!B \in \Gamma$. Upamana $!A \lor !B \in
\Gamma$, nanging $!A \notin \Gamma$ lan $!B \notin \Gamma$. Amarga
$\Gamma$ iku !!{complete}, $\lnot !A \in \Gamma$ lan $\lnot !B \in
\Gamma$. Miturut
\iftag{FOL}{%
  \tagrefs{prfSC/{fol:seq:ppr:prop:provability-lor},
    prfND/{fol:ntd:ppr:prop:provability-lor},
    prfAX/{fol:axd:ppr:prop:provability-lor},
    prfTab/{fol:tab:ppr:prop:provability-lor}}}{%
  \tagrefs{prfSC/{pl:seq:ppr:prop:provability-lor},
    prfND/{pl:ntd:ppr:prop:provability-lor},
    prfAX/{pl:axd:ppr:prop:provability-lor},
    prfTab/{pl:tab:ppr:prop:provability-lor}}},
butir (1), $\Gamma$ ora konsisten, kang dadi cengkah. Mula,
$!A \in \Gamma$ utawa $!B \in \Gamma$.

Kanggo arah kosok baline, upamana $!A \in \Gamma$ utawa $!B \in
\Gamma$. Miturut
\iftag{FOL}{%
  \tagrefs{prfSC/{fol:seq:ppr:prop:provability-lor},
    prfND/{fol:ntd:ppr:prop:provability-lor},
    prfAX/{fol:axd:ppr:prop:provability-lor},
    prfTab/{fol:tab:ppr:prop:provability-lor}}}{%
  \tagrefs{prfSC/{pl:seq:ppr:prop:provability-lor},
    prfND/{pl:ntd:ppr:prop:provability-lor},
    prfAX/{pl:axd:ppr:prop:provability-lor},
    prfTab/{pl:tab:ppr:prop:provability-lor}}}, butir (2),
$\Gamma \Proves !A \lor !B$. Miturut \olref{prop:ccs-prov-in},
$!A \lor !B \in \Gamma$, kaya kang dibutuhake.}}

\tagitem{defIf}{}{%
\iftag{probIf}{Gladhen.}{%
Kanggo arah maju, upamana $!A \lif !B \in \Gamma$, lan kanggo nggayuh
cengkah upamana $!A \in \Gamma$ lan $!B \notin \Gamma$. Miturut
hipotesis-hipotesis iki, $!A \lif !B \in \Gamma$ lan $!A \in
\Gamma$. Miturut
\iftag{FOL}{%
  \tagrefs{prfSC/{fol:seq:ppr:prop:provability-lif},
    prfND/{fol:ntd:ppr:prop:provability-lif},
    prfAX/{fol:axd:ppr:prop:provability-lif},
    prfTab/{fol:tab:ppr:prop:provability-lif}}}{%
  \tagrefs{prfSC/{pl:seq:ppr:prop:provability-lif},
    prfND/{pl:ntd:ppr:prop:provability-lif},
    prfAX/{pl:axd:ppr:prop:provability-lif},
    prfTab/{pl:tab:ppr:prop:provability-lif}}}, butir (1),
$\Gamma \Proves !B$. Nanging banjur, miturut
\olref{prop:ccs-prov-in}, $!B \in \Gamma$, kang cengkah karo hipotesis
yen $!B \notin \Gamma$.

Kanggo arah kosok baline, dhisik gatekna perkara $!A \notin \Gamma$.
Amarga $\Gamma$ iku !!{complete}, $\lnot !A \in \Gamma$. Miturut
\iftag{FOL}{%
  \tagrefs{prfSC/{fol:seq:ppr:prop:provability-lif},
    prfND/{fol:ntd:ppr:prop:provability-lif},
    prfAX/{fol:axd:ppr:prop:provability-lif},
    prfTab/{fol:tab:ppr:prop:provability-lif}}}{%
  \tagrefs{prfSC/{pl:seq:ppr:prop:provability-lif},
    prfND/{pl:ntd:ppr:prop:provability-lif},
    prfAX/{pl:axd:ppr:prop:provability-lif},
    prfTab/{pl:tab:ppr:prop:provability-lif}}}, butir (2),
$\Gamma \Proves !A \lif !B$. Miturut \olref{prop:ccs-prov-in} maneh,
kita entuk $!A \lif !B \in \Gamma$, kaya kang dibutuhake.

Saiki gatekna perkara $!B \in \Gamma$. Miturut
\iftag{FOL}{%
  \tagrefs{prfSC/{fol:seq:ppr:prop:provability-lif},
    prfND/{fol:ntd:ppr:prop:provability-lif},
    prfTab/{fol:tab:ppr:prop:provability-lif},
    prfAX/{fol:axd:ppr:prop:provability-lif}}}{%
  \tagrefs{prfSC/{pl:seq:ppr:prop:provability-lif},
    prfND/{pl:ntd:ppr:prop:provability-lif},
    prfTab/{pl:tab:ppr:prop:provability-lif},
    prfAX/{pl:axd:ppr:prop:provability-lif}}},
butir (2) maneh, $\Gamma \Proves !A \lif !B$. Miturut
\olref{prop:ccs-prov-in}, $!A \lif !B \in \Gamma$.}}
\end{enumerate}
\end{proof}

\tagprob{FOL}
\begin{prob}
Rampungna bukti \olref[fol][com][ccs]{prop:ccs}.
\end{prob}
\tagendprob

\tagprob{notFOL}
\begin{prob}
Rampungna bukti \olref[pl][com][ccs]{prop:ccs}.
\end{prob}
\tagendprob

\end{document}
